\documentclass[11pt,a4paper]{article}
\usepackage[utf8]{inputenc}
\usepackage[T1]{fontenc}
\usepackage{amsmath,amssymb}
\usepackage[margin=2.5cm]{geometry}
\usepackage{xcolor}

\setlength{\parindent}{0pt}
\setlength{\parskip}{6pt}

\begin{document}

\begin{center}
{\large Technical Neural Networks Assignment Sheet 5}\\[4pt]
{\bfseries Names: Mohamadhosein Sabaghi (50547423)}
\end{center}

\section*{Assignment 21}

\textbf{Idea.} Plain gradient descent uses the gradient's \emph{magnitude} directly as (part of) the
step size, $\Delta w_{ij}=-\eta\,\partial E/\partial w_{ij}$. But that magnitude is unreliable: with
saturating transfer functions (tanh, logistic) the gradient can be tiny in flat regions and large
elsewhere, so a single global $\eta$ is either too small for the flat parts or too large for the steep
parts. Resilient Propagation (RPROP) fixes this by using \emph{only the sign} of the gradient to decide
the update direction, and tracking a separate, adaptively-sized step $\Delta_{ij}$ for every single
weight, based purely on whether that weight's gradient kept pointing the same way or flipped.

\textbf{Step-size adaptation.} Let $g_{ij}(t):=\partial\,{}^pE/\partial w_{ij}$ at step $t$. Compare its
sign to the previous step:
\[
  \Delta_{ij}(t) =
  \begin{cases}
    \min(\eta^{+}\Delta_{ij}(t-1),\ \Delta_{\max}) & \text{if } g_{ij}(t-1)\,g_{ij}(t) > 0 \\
    \max(\eta^{-}\Delta_{ij}(t-1),\ \Delta_{\min}) & \text{if } g_{ij}(t-1)\,g_{ij}(t) < 0 \\
    \Delta_{ij}(t-1) & \text{if } g_{ij}(t-1)\,g_{ij}(t) = 0
  \end{cases}
\]
with $0<\eta^{-}<1<\eta^{+}$ (typically $\eta^+=1.2,\ \eta^-=0.5$).

\textbf{Weight update.} The step size $\Delta_{ij}$, not the gradient, determines how far to move:
\[
  \Delta w_{ij}(t) = -\operatorname{sign}\!\big(g_{ij}(t)\big)\cdot \Delta_{ij}(t), \qquad
  w_{ij}(t+1) = w_{ij}(t) + \Delta w_{ij}(t)
\]

\textbf{What each term means / motivation.}
\begin{itemize}
  \item $g_{ij}(t-1)\,g_{ij}(t)>0$: the gradient kept the same sign two steps in a row, i.e.\ we're
  still heading downhill in the same direction (e.g.\ on a long shallow slope) -- so $\Delta_{ij}$ is
  \emph{increased} by the acceleration factor $\eta^+$ to cross such regions faster.
  \item $g_{ij}(t-1)\,g_{ij}(t)<0$: the sign flipped, meaning the previous step overshot a local
  minimum along this weight's direction -- so $\Delta_{ij}$ is \emph{decreased} by $\eta^-$ to approach
  it more carefully and stop oscillating.
  \item $\Delta_{\min},\Delta_{\max}$: safety bounds so the per-weight step size can never shrink to
  exactly zero (permanently stuck) or grow without bound (diverging).
  \item Using $\operatorname{sign}(g_{ij})$ instead of $g_{ij}$ itself for the update direction is what
  decouples the step size from the gradient's magnitude, which is the whole point: a weight with a
  tiny but persistently-signed gradient still gets to move by a healthy, self-accelerating amount.
\end{itemize}

\textbf{Reference:} M.\ Riedmiller and H.\ Braun, \emph{``A direct adaptive method for faster
backpropagation learning: the RPROP algorithm''}, Proceedings of the IEEE International Conference on
Neural Networks (ICNN), 1993, pp.\ 586--591.

\section*{Assignment 22}

\textbf{Single-step (online) learning.} One pattern $p$ is presented, its $\Delta w_{ij}$ computed, and
the weights updated \emph{immediately}, before the next pattern is shown.

\textbf{Cumulative (batch) learning.} All $P$ patterns are presented first, their weight changes
${}^p\Delta w_{ij}$ accumulated, $\Delta w_{ij}=\sum_{p=1}^P {}^p\Delta w_{ij}$, and only then applied
once -- exact gradient descent on the total error $G=\sum_p {}^pE$.

\begin{itemize}
  \item \textbf{Single-step:} \emph{+} many more updates per epoch, often converges faster in
  wall-clock time; the noise from updating on one pattern at a time can help escape shallow local
  minima; no extra memory needed to accumulate $\Delta w$; works for streaming data. \emph{--} the
  update direction is noisy and zig-zags (it follows one pattern's gradient, not $G$'s true gradient);
  the result depends on the order patterns are presented in.
  \item \textbf{Cumulative:} \emph{+} follows the true gradient of $G$, so convergence is smooth and
  order-independent; easier to analyze with standard gradient-descent theory. \emph{--} needs a full
  pass over all patterns before a single update (slower per update, more memory); with no stochastic
  noise it's more likely to get stuck exactly in a local minimum or saddle point; not usable in an
  online/streaming setting where the full pattern set isn't available upfront.
\end{itemize}

\section*{Assignment 23}

A schedule like $\eta(t)=\eta_0/t$ shrinks very quickly and keeps shrinking forever, regardless of how
close training actually is to a good solution: after only a few steps $\eta$ is already a small fraction
of $\eta_0$, no matter what the error surface still looks like at that point. The problem is that this
decay is tied purely to elapsed time $t$, not to the network's actual progress. If the weights still
happen to be somewhere in the middle of a large or complicated error surface once $\eta$ has already
shrunk to a tiny value, the steps become too small to make further real progress, but training hasn't
reached a minimum yet -- it just grinds to a halt, stuck exactly where the clock happened to catch up
with it. The safer practical alternative is to reduce $\eta$ by a fixed factor (e.g.\ $\eta\leftarrow
0.1\,\eta$) only at specific, deliberately chosen moments, rather than continuously via a fixed formula
in $t$. (This is also why RPROP from assignment 21 is a better idea: its step size only shrinks when
there is actual evidence of overshooting -- a gradient sign flip -- not just because time has passed.)

\section*{Assignment 24}

Identity transfer function means every layer's output equals its net input: $f(z)=z$. Write the MLP in
vector-matrix form, omitting the BIAS as permitted.

\textbf{3 layers (one hidden layer).} Input $\underline X\in\mathbb R^N$, weight matrix
$\underline U$ ($H\times N$) into the hidden layer, weight matrix $\underline V$ ($M\times H$) into the
output layer:
\[
  \underline{\mathrm{out}}_{\text{hidden}} = \underline U\,\underline X, \qquad
  \underline Y = \underline V\,\underline{\mathrm{out}}_{\text{hidden}} = \underline V\,(\underline U\,\underline X) = (\underline V\,\underline U)\,\underline X
\]
Matrix multiplication is associative, so $\underline V\,\underline U$ is just one more $M\times N$
matrix -- call it $\underline W:=\underline V\,\underline U$. Then
\[
  \boxed{\underline Y = \underline W\,\underline X}
\]
which is \emph{exactly} the mapping of a 2-layer Perceptron with weight matrix $\underline W$ and
identity transfer function: the 3-layer MLP and this Perceptron compute the identical function
$\underline Y(\underline X)$ for every input, with $\underline W=\underline V\,\underline U$.

\textbf{More than 3 layers (induction).} For $L$ weight matrices $\underline{W}_1,\dots,\underline{W}_L$
(so $L+1$ layers total),
\[
  \underline Y = \underline{W}_L\big(\underline{W}_{L-1}(\cdots(\underline{W}_1\underline X)\cdots)\big)
  = (\underline{W}_L\,\underline{W}_{L-1}\cdots\underline{W}_1)\,\underline X
\]
again by associativity -- any chain of matrix multiplications collapses into a single matrix
$\underline W := \underline{W}_L\cdots\underline{W}_1$, so $\underline Y=\underline W\underline X$ once
more. (Formally: assume an MLP with $L$ weight matrices collapses to a single $\underline W_L'$; adding
one more layer with matrix $\underline W_{L+1}$ gives
$\underline Y=\underline W_{L+1}(\underline W_L'\underline X)=(\underline W_{L+1}\underline W_L')\underline X$,
still a single matrix -- so the claim holds for every depth by induction.)

So no matter how many layers an all-identity MLP has, it computes a purely linear map
$\underline X\mapsto\underline W\underline X$, exactly reproducible by a single 2-layer Perceptron --
depth adds no extra expressive power without a nonlinear transfer function. $\blacksquare$

\section*{Assignment 25}

\textbf{Overfitting} (``learning by heart'') happens when a network ends up fitting the specific noise
and idiosyncrasies of the training patterns instead of the underlying regularity behind them -- it then
performs very well on the training set but poorly on new, unseen data (checked against a validation set
not used for training). It tends to occur when the network has too many free parameters (weights)
relative to how much training data and how simple the true underlying task actually is.

\textbf{Strategy to avoid it: weight decay.} Extend the error function with a term that penalizes large
weights,
\[
  {}^pE^{**} = {}^pE + \tfrac12\beta\sum_{i,j}w_{ij}^2,
\]
exactly the one derived in assignment 19 (sheet 3), giving the update
$\Delta w_{ij}=\eta\,\delta_j\,\mathrm{out}_i-\eta\beta w_{ij}$. The extra $-\eta\beta w_{ij}$ term
continuously pulls every weight toward $0$; weights that are driven to (near) zero can no longer
contribute to memorizing individual training patterns, which reduces the network's effective capacity
and pushes it toward the smoother, simpler mapping that generalizes better.

\section*{Assignment 26}

\textbf{AI use:} I used Claude (Anthropic) to work through this part, as the sheet requires. I gave it
the assignment text and asked for an original example illustrating generalization in neural networks.
The example below was usable as given, so I didn't need to re-prompt or discard anything.

\textbf{Generalization} is a network's ability to produce good outputs for inputs it was \emph{not}
trained on, by having learned the true underlying regularity rather than memorizing the training set.

\textbf{Example.} Train a small MLP to decide whether a 2D point $(x_1,x_2)$ lies inside or outside a
circle of radius $r$ centered at the origin, using only $20$ labeled training points scattered near the
boundary. A network with too few hidden neurons might draw a decision boundary that is too simple (e.g.\
close to a straight line) and misclassify many points even in training -- underfitting. A network with
far too many hidden neurons might instead draw a highly wiggly boundary that perfectly separates exactly
those $20$ training points, including the noise in their exact positions -- overfitting: on the training
set the error is $0$, but a new test point sitting right next to a training point can end up on the
wrong side of one of those wiggles. A network with an appropriately sized hidden layer instead settles
on a boundary close to the true circle, correctly classifying new, previously unseen $(x_1,x_2)$ points
near the boundary that it never saw during training -- that is good generalization.

\end{document}
